% Generated by build/build.R from data/schedule.yml. Do not edit by hand.

\subsection*{Class 1 -- Overview}
\begin{itemize}
	\item[$\square$] Syllabus
\end{itemize}

\subsection*{Class 2 -- Introduction to environmental economics}
\begin{itemize}
	\item[$\square$] Don Fullerton and Robert Stavins (1998). ``How economists see the environment.'' \textit{Nature}. \href{https://scholar.harvard.edu/files/stavins/files/how_economists_see_the_environment.pdf}{[link]}
	\item[$\square$] Kenneth J. Arrow et al. (1996). ``Is there a role for benefit-cost analysis in environmental, health, and safety regulation?'' \textit{Science}.
	\item[$\square$] John V. Krutilla (1967). ``Conservation reconsidered.'' \textit{American Economic Review}.
\end{itemize}

\subsection*{Class 3 -- Theory of externalities}
\begin{itemize}
	\item[$\square$] P\&R (Chapters 1--2)
	\item[$\square$] Paul A. Samuelson (1954). ``The pure theory of public expenditure.'' \textit{Review of Economics and Statistics}.
\end{itemize}

\subsection*{Class 4 -- Theory of environmental policy}
\begin{itemize}
	\item[$\square$] P\&R (Chapter 3)
	\item[$\square$] Ronald Coase (1960). ``The problem of social cost.'' \textit{Journal of Law and Economics}.
\end{itemize}

\subsection*{Class 5 -- Instrument choice I}
\begin{itemize}
	\item[$\square$] P\&R (Chapter 4)
	\item[$\square$] Martin Weitzman (1974). ``Prices versus quantities.'' \textit{Review of Economic Studies}.
\end{itemize}

\subsection*{Class 6 -- Instrument choice II}
\begin{itemize}
	\item[$\square$] P\&R (Chapter 5)
	\item[$\square$] Lawrence Goulder and Ian Parry (2008). ``Instrument choice in environmental policy.'' \textit{Review of Environmental Economics and Policy}.
\end{itemize}

\subsection*{Class 7 -- Instrument choice III}
\begin{itemize}
	\item[$\square$] P\&R (Chapter 5, continued)
	\item[$\square$] Thomas H. Tietenberg (1990). ``Economic instruments for environmental regulation.'' \textit{Oxford Review of Economic Policy}.
\end{itemize}

\subsection*{Class 8 -- Instrument choice IV}
\begin{itemize}
	\item[$\square$] P\&R (Chapter 7)
	\item[$\square$] Lawrence Goulder, Ian Parry, Roberton C. Williams III, and Dallas Burtraw (1999). ``The cost-effectiveness of alternative instruments for environmental protection in a second-best setting.'' \textit{Journal of Public Economics}.
	\item[$\square$] Sarah E. West and Roberton C. Williams III (2007). ``Optimal taxation and cross-price effects on labor supply: Estimates of the optimal gas tax.'' \textit{Journal of Public Economics}.
\end{itemize}

\subsection*{Class 9 -- Welfare analysis}
\begin{itemize}
	\item[$\square$] P\&R (Chapter 14)
	\item[$\square$] Robert D. Willig (1976). ``Consumer's surplus without apology.'' \textit{American Economic Review}.
\end{itemize}

\subsection*{Class 10 -- Valuation I -- revealed preference}
\begin{itemize}
	\item[$\square$] P\&R (Chapter 15)
	\item[$\square$] W. Michael Hanemann (1991). ``Willingness to pay and willingness to accept: How much can they differ?'' \textit{American Economic Review}.
	\item[$\square$] Jason F. Shogren, Seung Y. Shin, Dermot J. Hayes, and James B. Kliebenstein (1994). ``Resolving differences in willingness to pay and willingness to accept.'' \textit{American Economic Review}.
\end{itemize}

\subsection*{Class 11 -- Valuation II -- property value models}
\begin{itemize}
	\item[$\square$] P\&R (Chapter 18, focusing on 18.1 and 18.2)
	\item[$\square$] Kenneth Y. Chay and Michael Greenstone (2005). ``Does air quality matter? Evidence from the housing market.'' \textit{Journal of Political Economy}.
	\item[$\square$] H. Spencer Banzhaf and Randall P. Walsh (2008). ``Do people vote with their feet? An empirical test of Tiebout.'' \textit{American Economic Review}.
	\item[$\square$] Peter Christensen and Christopher Timmins (2022). ``Sorting or steering: The effects of housing discrimination on neighborhood choice.'' \textit{Journal of Political Economy}.
\end{itemize}

\subsection*{Class 12 -- Attend SWEEEP in lieu of class}
\noindent Readings to be announced.

\subsection*{Class 13 -- Guest lecture -- hedonics and VSL}
\noindent Readings to be announced.

\subsection*{Class 14 -- Marginal cost of public funds}
\begin{itemize}
	\item[$\square$] Amy Finkelstein and Nathaniel Hendren (2020). ``Welfare analysis meets causal inference.'' \textit{Journal of Economic Perspectives}.
	\item[$\square$] Robert W. Hahn, Nathaniel Hendren, Robert D. Metcalfe, and Ben Sprung-Keyser (2024). ``A welfare analysis of policies impacting climate change.'' \textit{NBER Working Paper 32728}. \href{https://www.nber.org/papers/w32728}{[link]}
\end{itemize}

\subsection*{Class 15 -- Midterm (in class)}

\subsection*{Class 16 -- TBD}
\noindent Readings to be announced.

\subsection*{Class 17 -- Topics: Utility pricing}
\begin{itemize}
	\item[$\square$] Casey J. Wichman (2024). ``Efficiency, equity, and cost-recovery trade-offs in municipal water pricing.'' \textit{RFF Working Paper}. \hfill \textbf{[response due]}
	\item[$\square$] Arik Levinson and Emilson Silva (2022). ``The electric Gini: Income redistribution through energy prices.'' \textit{American Economic Journal: Economic Policy}.
	\item[$\square$] Koichiro Ito (2014). ``Do consumers respond to marginal or average price? Evidence from nonlinear electricity pricing.'' \textit{American Economic Review}.
\end{itemize}
\noindent \textit{Discussion led by Prof. Wichman.}

\subsection*{Class 18 -- Topics: Water}
\begin{itemize}
	\item[$\square$] David A. Keiser, Bhashkar Mazumder, David Molitor, and Joseph S. Shapiro (2023). ``Water works: Causes and consequences of safe drinking water in America.'' \textit{Working Paper}.
	\item[$\square$] W. Michael Hanemann (2006). ``The economic conception of water.'' \textit{Water Crisis: Myth or Reality?}.
	\item[$\square$] Sheila M. Olmstead (2010). ``The economics of managing scarce water resources.'' \textit{Review of Environmental Economics and Policy}.
	\item[$\square$] Andrew B. Ayres, Kyle C. Meng, and Andrew J. Plantinga (2021). ``Do environmental markets improve on open access? Evidence from California groundwater rights.'' \textit{Journal of Political Economy}.
	\item[$\square$] Charles A. Taylor and Hannah Druckenmiller (2022). ``Wetlands, flooding, and the Clean Water Act.'' \textit{American Economic Review}.
\end{itemize}

\subsection*{Class 19 -- Topics: Electricity externalities and pricing}
\begin{itemize}
	\item[$\square$] Stephen P. Holland, Erin T. Mansur, Nicholas Z. Muller, and Andrew J. Yates (2016). ``Are there environmental benefits from driving electric vehicles? The importance of local factors.'' \textit{American Economic Review}.
	\item[$\square$] Severin Borenstein and James B. Bushnell (2022). ``Do two electricity pricing wrongs make a right? Cost recovery, externalities, and efficiency.'' \textit{American Economic Journal: Economic Policy}.
\end{itemize}

\subsection*{Class 20 -- Topics: Energy and smart technologies}
\begin{itemize}
	\item[$\square$] Joshua Blonz, Karen Palmer, Casey J. Wichman, and Derek C. Wietelman (2025). ``Smart thermostats, automation, and time-varying prices.'' \textit{American Economic Journal: Applied Economics}.
	\item[$\square$] Megan R. Bailey, David P. Brown, Blake C. Shaffer, and Frank A. Wolak (2025). ``Take the load off: Time and technology as determinants of electricity demand response.'' \textit{NBER Working Paper 33836}. \href{https://www.nber.org/papers/w33836}{[link]}
	\item[$\square$] Alec Brandon, Christopher M. Clapp, John A. List, Robert D. Metcalfe, and Michael Price (2022). ``The human perils of scaling smart technologies: Evidence from field experiments.'' \textit{NBER Working Paper 30482}. \href{https://www.nber.org/papers/w30482}{[link]}
\end{itemize}

\subsection*{Class 21 -- Topics: Climate change}
\begin{itemize}
	\item[$\square$] Tamma Carleton et al. (2022). ``Valuing the global mortality consequences of climate change accounting for adaptation costs and benefits.'' \textit{Quarterly Journal of Economics}.
	\item[$\square$] Alan Barreca, Karen Clay, Olivier Deschenes, Michael Greenstone, and Joseph S. Shapiro (2016). ``Adapting to climate change: The remarkable decline in the US temperature-mortality relationship over the twentieth century.'' \textit{Journal of Political Economy}.
	\item[$\square$] Marshall Burke, Solomon M. Hsiang, and Edward Miguel (2015). ``Global non-linear effect of temperature on economic production.'' \textit{Nature}.
	\item[$\square$] Kevin Rennert et al. (2022). ``Comprehensive evidence implies a higher social cost of CO2.'' \textit{Nature}.
	\item[$\square$] Tamma Carleton and Michael Greenstone (2022). ``A guide to updating the US government's social cost of carbon.'' \textit{Review of Environmental Economics and Policy}.
	\item[$\square$] Robert S. Pindyck (2013). ``Climate change policy: What do the models tell us?'' \textit{Journal of Economic Literature}.
	\item[$\square$] Wolfram Schlenker and Michael J. Roberts (2009). ``Nonlinear temperature effects indicate severe damages to US crop yields under climate change.'' \textit{Proceedings of the National Academy of Sciences}.
\end{itemize}

\subsection*{Class 22 -- Topics: Conservation}
\begin{itemize}
	\item[$\square$] Eyal Frank and Anant Sudarshan (2024). ``The social costs of keystone species collapse: Evidence from the decline of vultures in India.'' \textit{American Economic Review}.
\end{itemize}

\subsection*{Class 23 -- Topics: Forecasts}
\begin{itemize}
	\item[$\square$] Renato Molina and Ivan Rudik (2024). ``The social value of hurricane forecasts.'' \textit{NBER Working Paper 32548}. \href{https://www.nber.org/papers/w32548}{[link]}
\end{itemize}

\subsection*{Class 24 -- Topics: Domestic climate policy}
\begin{itemize}
	\item[$\square$] Hunt Allcott, Reigner Kane, Maximilian S. Maydanchik, Joseph S. Shapiro, and Felix Tintelnot (2024). ``The effects of Buy American: Electric vehicles and the Inflation Reduction Act.'' \textit{NBER Working Paper 33032}. \href{https://www.nber.org/papers/w33032}{[link]}
	\item[$\square$] John Bistline, Kimberly A. Clausing, Neil R. Mehrotra, James H. Stock, and Catherine Wolfram (2025). ``Climate policy reform options in 2025.'' \textit{Environmental and Energy Policy and the Economy}.
\end{itemize}

\subsection*{Class 25 -- Proposal presentations}

\subsection*{Class 26 -- Proposal presentations}

\subsection*{Class 27 -- Proposal presentations}

\subsection*{Class 28 -- Proposal presentations}
